% =========================================================================
% TEMPLATE TUGAS AKHIR PROGRAM STUDI SARJANA MATEMATIKA FMIPA ITB
% Sesuai Pedoman Penulisan Tugas Akhir - Januari 2017
% =========================================================================

\documentclass[12pt,a4paper,oneside]{book}

% -------------------------------------------------------------------------
% PAKET-PAKET UTAMA
% -------------------------------------------------------------------------
\usepackage[utf8]{inputenc}
\usepackage[T1]{fontenc}
\usepackage{times}            % Huruf Times New Roman (Font 12)
\usepackage[indonesian]{babel}
\usepackage{amsmath,amssymb,amsfonts,amsthm}
\usepackage{bm}
\usepackage{graphicx}
\usepackage{setspace}         % Pengaturan spasi baris
\usepackage{booktabs}
\usepackage{tabularx}
\usepackage{caption}
\usepackage{tocbibind}        % Memasukkan Daftar Pustaka ke Daftar Isi
\usepackage{hyperref}

\hypersetup{
    colorlinks=false,         % Cetakan resmi hitam-putih
    pdfborder={0 0 0},
    pdfencoding=auto,
    psdextra
}

% -------------------------------------------------------------------------
% GEOMETRI HALAMAN (Bab III.2 & III.4 Pedoman ITB 2017)
% 1 sisi (tebal < 1.5 cm): kiri = 4 cm, atas = 3 cm, kanan = 3 cm, bawah = 3 cm
% Nomor halaman di tengah, 1.5 cm dari tepi bawah kertas (footskip)
% -------------------------------------------------------------------------
\usepackage{geometry}
\geometry{
    a4paper,
    left=40mm,
    top=30mm,
    right=30mm,
    bottom=30mm,
    footskip=15mm
}

% Jarak baris 1.5 spasi (Bab III.2.3)
\onehalfspacing

% Penomoran halaman di bawah tengah
\pagestyle{plain}

% Kedalaman Daftar Isi: Hanya sampai tingkat Bab (tanpa subbab)
\setcounter{tocdepth}{0}

% -------------------------------------------------------------------------
% PENGATURAN JUDUL BAB & FORMAT TEKS
% Sesuai Bab V.1: Judul bab 14pt tebal, simetris/centered
% -------------------------------------------------------------------------
\usepackage{titlesec}
\titleformat{\chapter}[display]
  {\normalfont\fontsize{14}{17}\bfseries\centering}
  {\MakeUppercase{\chaptertitlename\ \Roman{chapter}}}
  {0.5ex}
  {\fontsize{14}{17}\bfseries}
\titlespacing*{\chapter}{0pt}{-20pt}{25pt}

\titleformat{\section}
  {\normalfont\fontsize{12}{15}\bfseries}
  {\thesection}
  {1em}
  {}

\titleformat{\subsection}
  {\normalfont\fontsize{12}{15}\bfseries}
  {\thesubsection}
  {1em}
  {}

% Format penomoran subbab: II.1, II.2 (Roman.Arab)
\renewcommand{\thechapter}{\Roman{chapter}}
\renewcommand{\thesection}{\thechapter.\arabic{section}}
\renewcommand{\thesubsection}{\thesection.\arabic{subsection}}

% Format caption tabel & gambar: Nomor bab Arab diikuti nomor urut (contoh: Tabel 2.1)
\renewcommand{\thetable}{\arabic{chapter}.\arabic{table}}
\renewcommand{\thefigure}{\arabic{chapter}.\arabic{figure}}

% Penyesuaian kata bahasa Indonesia
\addto\captionsindonesian{%
  \renewcommand{\chaptername}{Bab}%
  \renewcommand{\contentsname}{DAFTAR ISI}%
  \renewcommand{\listfigurename}{DAFTAR GAMBAR DAN ILUSTRASI}%
  \renewcommand{\listtablename}{DAFTAR TABEL}%
  \renewcommand{\bibname}{DAFTAR PUSTAKA}%
  \renewcommand{\appendixname}{LAMPIRAN}%
}

% Lingkungan teorema matematika
\newtheorem{teorema}{Teorema}[chapter]
\newtheorem{definisi}{Definisi}[chapter]
\newtheorem{lemma}{Lemma}[chapter]
\newtheorem{contoh}{Contoh}[chapter]

% =========================================================================
% DOKUMEN UTAMA
% =========================================================================
\begin{document}

% -------------------------------------------------------------------------
% 1. SAMPUL LUAR / SAMPUL DALAM (Lampiran A)
% Tanpa nomor halaman
% -------------------------------------------------------------------------
\begin{titlepage}
    \centering
    \vspace*{0.5cm}
    {\fontsize{14}{17}\bfseries SEGMENTASI GAMBAR DENGAN DNN DENGAN BOBOT BERDISTRIBUSI NORMAL\par}
    
    \vspace{2.0cm}
    {\fontsize{14}{17}\bfseries TUGAS AKHIR\par}
    \vspace{0.6cm}
    {\fontsize{12}{15} Karya tulis sebagai salah satu syarat\\
    untuk memperoleh gelar Sarjana dari\\
    Institut Teknologi Bandung\par}
    
    \vspace{1.5cm}
    {\fontsize{12}{15}\bfseries Oleh\par}
    \vspace{0.3cm}
    {\fontsize{14}{17}\bfseries GULIT RADIAN WIYARNO\par}
    \vspace{0.2cm}
    {\fontsize{14}{17}\bfseries NIM: 10123033\par}
    \vspace{0.2cm}
    {\fontsize{14}{17}\bfseries (Program Studi Sarjana Matematika)\par}
    
    \vspace{1.5cm}
    \includegraphics[width=3.0cm]{figures/logo-itb-hitam.jpg}\par
    
    \vfill
    {\fontsize{14}{17}\bfseries INSTITUT TEKNOLOGI BANDUNG\par}
    \vspace{0.2cm}
    {\fontsize{14}{17}\bfseries Juni 2027\par}
    \vspace*{0.5cm}
\end{titlepage}

% -------------------------------------------------------------------------
% BAGIAN PERSIAPAN TUGAS AKHIR (Bab II.4 & III.3)
% Penomoran halaman angka Romawi kecil (i, ii, iii, ...)
% -------------------------------------------------------------------------
\frontmatter
\pagenumbering{roman}

% -------------------------------------------------------------------------
% 2. HALAMAN PENGESAHAN (Lampiran D)
% -------------------------------------------------------------------------
\newpage
\thispagestyle{plain}
\begin{center}
    \vspace*{0.5cm}
    {\fontsize{14}{17}\bfseries SEGMENTASI GAMBAR DENGAN DNN DENGAN BOBOT BERDISTRIBUSI NORMAL\par}
    
    \vspace{2.5cm}
    {\fontsize{12}{15} Oleh\par}
    \vspace{0.3cm}
    {\fontsize{12}{15}\bfseries Gulit Radian Wiyarno\par}
    \vspace{0.2cm}
    {\fontsize{12}{15}\bfseries NIM: 10123033\par}
    \vspace{0.2cm}
    {\fontsize{12}{15}(Program Studi Sarjana Matematika)\par}
    
    \vspace{0.8cm}
    {\fontsize{12}{15} Institut Teknologi Bandung\par}
    
    \vspace{1.5cm}
    {\fontsize{12}{15} Menyetujui\\
    Tim Pembimbing\par}
    
    \vspace{0.4cm}
    {\fontsize{12}{15} Tanggal .....................................\par}
    
    \vspace{1.5cm}
    \noindent
    \begin{minipage}[t]{0.48\textwidth}
        \centering
        {\fontsize{12}{15} Pembimbing Pertama\par}
        \vspace{2.0cm}
        \underline{\hspace{5.5cm}}\par
        \vspace{0.2cm}
        {\fontsize{11}{14}\bfseries Prof. Sapto Wahyu Indratno, S.Si., M.Sc., Ph.D.\par}
    \end{minipage}%
    \hfill
    \begin{minipage}[t]{0.48\textwidth}
        \centering
        {\fontsize{12}{15} Pembimbing Kedua\par}
        \vspace{2.0cm}
        \underline{\hspace{5.5cm}}\par
        \vspace{0.2cm}
        {\fontsize{11}{14}\bfseries Maharani Ahsani Ummi, S.Si., M.Si., M.Sc., Ph.D.\par}
    \end{minipage}\par
\end{center}

% -------------------------------------------------------------------------
% 3. HALAMAN PERUNTUKAN / PERSEMBAHAN (Opsional, Bab IV.3 - dinonaktifkan)
% -------------------------------------------------------------------------
% \newpage
% \thispagestyle{plain}
% \vspace*{6cm}
% \begin{flushright}
%     \textit{Kupersembahkan karya tulis ini untuk...}
% \end{flushright}

% -------------------------------------------------------------------------
% 4. ABSTRAK (BAHASA INDONESIA) (Bab II.2 & Lampiran B)
% -------------------------------------------------------------------------
\newpage
\thispagestyle{plain}
\begin{center}
    {\fontsize{14}{17}\bfseries ABSTRAK\par}
    \vspace{1.0cm}
    {\fontsize{14}{17}\bfseries SEGMENTASI GAMBAR DENGAN DNN DENGAN BOBOT BERDISTRIBUSI NORMAL\par}
    \vspace{0.8cm}
    {\fontsize{12}{15} Oleh\par}
    \vspace{0.2cm}
    {\fontsize{14}{17}\bfseries Gulit Radian Wiyarno\par}
    \vspace{0.2cm}
    {\fontsize{14}{17}\bfseries NIM: 10123033\par}
    \vspace{0.2cm}
    {\fontsize{14}{17}\bfseries (Program Studi Sarjana Matematika)\par}
\end{center}

\vspace{1.0cm}
\noindent
Abstrak Tugas Akhir memuat secara komprehensif permasalahan penelitian yang akan dilaksanakan dengan menjelaskan masalah ilmiah yang telah diteliti. Abstrak Tugas Akhir mencakup latar belakang, permasalahan yang dikaji, tahapan, tujuan, metode yang digunakan, ulasan singkat, serta penjelasan hasil dan kesimpulan yang diperoleh dalam penelitian. Abstrak terdiri atas 500 -- 800 kata dan di dalam abstrak tidak boleh ada rujukan referensi (tidak boleh ada hasil kajian dari referensi).

\vspace{0.8cm}
\noindent
\textbf{Kata kunci:} kata kunci 1, kata kunci 2, kata kunci 3.

% -------------------------------------------------------------------------
% 5. ABSTRACT (BAHASA INGGRIS) (Bab II.3)
% -------------------------------------------------------------------------
\newpage
\thispagestyle{plain}
\begin{center}
    {\fontsize{14}{17}\bfseries ABSTRACT\par}
    \vspace{1.0cm}
    {\fontsize{14}{17}\bfseries IMAGE SEGMENTATION USING DNN WITH NORMALLY DISTRIBUTED WEIGHTS\par}
    \vspace{0.8cm}
    {\fontsize{12}{15} By\par}
    \vspace{0.2cm}
    {\fontsize{14}{17}\bfseries Gulit Radian Wiyarno\par}
    \vspace{0.2cm}
    {\fontsize{14}{17}\bfseries NIM: 10123033\par}
    \vspace{0.2cm}
    {\fontsize{14}{17}\bfseries (Undergraduate Program in Mathematics)\par}
\end{center}

\vspace{1.0cm}
\noindent
\textit{The abstract comprehensively contains the research problems, objectives, methods, stages of research, concise discussion, and results/conclusions obtained. Written in 500 -- 800 words without literature citations.}

\vspace{0.8cm}
\noindent
\textit{\textbf{Keywords:} keyword 1, keyword 2, keyword 3.}

% -------------------------------------------------------------------------
% 6. KATA PENGANTAR (Bab IV.4)
% -------------------------------------------------------------------------
\newpage
\thispagestyle{plain}
\begin{center}
    {\fontsize{14}{17}\bfseries KATA PENGANTAR\par}
\end{center}
\vspace{0.8cm}

Puji dan syukur penulis panjatkan ke hadirat Tuhan Yang Maha Esa atas segala rahmat dan karunia-Nya sehingga Tugas Akhir ini dapat diselesaikan. Tugas Akhir ini disusun sebagai salah satu syarat untuk memperoleh gelar Sarjana dari Institut Teknologi Bandung.

% Tuliskan ucapan terima kasih kepada pihak-pihak terkait di sini saat naskah siap.

\vspace{1.5cm}
\begin{flushright}
    Bandung, Juni 2027\\[1.8cm]
    \underline{\hspace{5cm}}\\[0.2cm]
    Gulit Radian Wiyarno\\
    NIM: 10123033
\end{flushright}

% -------------------------------------------------------------------------
% 7. DAFTAR ISI, DAFTAR GAMBAR, DAFTAR TABEL (Bab IV.5 & IV.6)
% Jarak baris satu spasi (Bab IV.5)
% -------------------------------------------------------------------------
\newpage
\begin{singlespace}
    \tableofcontents
\end{singlespace}

\newpage
\begin{singlespace}
    \listoffigures
\end{singlespace}

\newpage
\begin{singlespace}
    \listoftables
\end{singlespace}

% -------------------------------------------------------------------------
% TUBUH UTAMA TUGAS AKHIR (Bab II.5, III.3, & V)
% Penomoran angka Arab (1, 2, 3, ...)
% -------------------------------------------------------------------------
\mainmatter
\pagenumbering{arabic}

\chapter{Pendahuluan}
\label{chap:pendahuluan}

\section{Latar Belakang}
\label{sec:latar_belakang}
% Tuliskan latar belakang penelitian di sini.

\section{Rumusan Masalah}
\label{sec:rumusan_masalah}
% Tuliskan rumusan masalah / pertanyaan penelitian di sini.

\section{Tujuan Penelitian}
\label{sec:tujuan}
% Tuliskan tujuan penelitian di sini.

\section{Batasan Masalah}
\label{sec:batasan_masalah}
% Tuliskan batasan atau asumsi yang digunakan di sini.

\section{Metode Penelitian}
\label{sec:metode_penelitian}
% Tuliskan ringkasan metode/pendekatan penelitian di sini.

\section{Sistematika Penulisan}
\label{sec:sistematika}
% Tuliskan outline/sistematika penulisan tugas akhir di sini.

\chapter{Landasan Teori}
\label{chap:landasan_teori}

\begin{enumerate}
    \item Gaussian Process
    \item Covariance Functions (Fungsi Kernel)
    \item Operasi Kombinasi Kernel
    \item Gaussian Process Regression (GPR)
    \item Sampling Jalur Fungsi Gaussian Process (Dekomposisi Cholesky)
    \item Marginal Likelihood dan Log Marginal Likelihood (LML)
    \item Optimasi Hyperparameter Berbasis Gradien
    \item Deep Gaussian Process (DGP)
    \item Deep Neural Network (DNN)
    \item Forward Propagation (Tingkat Neuron, Sampel, dan Batch)
    \item Backward Propagation (Gradient Descent dan Rantai Turunan Matriks)
    \item Bobot Berdistribusi Normal (Variational Weights)
    \item Reparameterization Trick
    \item Estimasi Ketidakpastian dan Interval Persentil
\end{enumerate}

\chapter{Model}
\label{chap:model}

\section{Tahapan Penelitian}
\label{sec:tahapan_penelitian}
% Uraikan alur kerja atau tahapan pelaksanaan penelitian.

\section{Deskripsi dan Pengolahan Data}
\label{sec:data}
% Uraikan dataset, sumber data, dan teknik preprocessing yang digunakan.

\section{Formulasi Model}
\label{sec:formulasi_model}
% Uraikan detail matematis atau arsitektur model yang dibangun.

\chapter{Hasil dan Pembahasan}
\label{chap:hasil_pembahasan}

\section{Hasil Eksperimen}
\label{sec:hasil_eksperimen}
% Paparkan hasil eksperimen, pengamatan, atau pengujian.

\section{Analisis dan Pembahasan}
\label{sec:analisis_pembahasan}
% Bahas analisis dari temuan atau hasil eksperimen secara mendalam.

\chapter{Kesimpulan dan Saran}
\label{chap:kesimpulan}

\section{Kesimpulan}
\label{sec:kesimpulan}
% Tuliskan rincian kesimpulan yang menjawab tujuan penelitian.

\section{Saran}
\label{sec:saran}
% Tuliskan saran untuk penelitian lanjutan atau implikasi praktis.


% -------------------------------------------------------------------------
% DAFTAR PUSTAKA (Bab II.6 & Lampiran E)
% Menggunakan satu spasi per entri
% -------------------------------------------------------------------------
\newpage
\begin{singlespace}
    \bibliographystyle{apalike}
    \bibliography{references}
\end{singlespace}

% -------------------------------------------------------------------------
% LAMPIRAN (Bab II.7)
% Didahului halaman pemisah bertuliskan LAMPIRAN di tengah halaman
% Nomor halaman kelanjutan angka Arab dari tubuh utama
% -------------------------------------------------------------------------
\newpage
\thispagestyle{plain}
\vspace*{\fill}
\begin{center}
    {\fontsize{14}{17}\bfseries LAMPIRAN}
\end{center}
\vspace*{\fill}

\appendix
% Penomoran lampiran: Lampiran A, B, C, dst.
\titleformat{\chapter}[display]
  {\normalfont\fontsize{14}{17}\bfseries\centering}
  {\MakeUppercase{\appendixname\ \Alph{chapter}}}
  {0.5ex}
  {\fontsize{14}{17}\bfseries}

\renewcommand{\thechapter}{\Alph{chapter}}
\renewcommand{\thesection}{\thechapter.\arabic{section}}
\renewcommand{\thetable}{\Alph{chapter}.\arabic{table}}
\renewcommand{\thefigure}{\Alph{chapter}.\arabic{figure}}

\chapter{Data Tambahan dan Penurunan Rumus}
\label{app:data_tambahan}

% Tuliskan materi lampiran A di sini (penurunan rumus, data mentah, kode, dsb.).


\end{document}
